\documentclass[aip,adv,reprint,floatfix]{revtex4-2}

\usepackage{amsmath,amssymb}
\usepackage{graphicx}

\begin{document}

\title{Geometric Photonic Quantum Error Correction via SEFI, GWFM, and DEFI}

\author{Jason Dutton}
\affiliation{Castlewood, Virginia, USA}

\date{13 September 2026}

\begin{abstract}
Photonic quantum error correction has traditionally relied on algebraic and phase-space constructions such as
cat, binomial, and GKP codes.\cite{DuttonSEFI2026,DuttonElectronField2026} These frameworks provide strong protection against loss and dephasing but
do not explicitly incorporate the geometric constraints that govern photonic mode evolution. The SEFI field
architecture\cite{DuttonSEFI2026} and the Geometric Worldline Field Model (GWFM)\cite{DuttonElectronField2026} introduce stability surfaces and worldline
geometry as physically meaningful structures for describing photonic modes, errors, and correction. In this
work, Dynamic Entity Field Integration (DEFI) is introduced as the unifying field law that binds SEFI identity
and GWFM geometry into a single stability surface. This DEFI-centric framework provides a geometric
foundation for photonic QEC, enabling curvature-based stabilization, geometric filtering, and sovereignty
restoration across diverse photonic platforms. Figure~\ref{fig:stability_surface} illustrates the DEFI stability
surface that anchors this unified geometric interpretation.
\end{abstract}

\maketitle

\section{BACKGROUND AND MOTIVATION}

Photonic quantum error correction has traditionally relied on algebraic and phase-space constructions such as
cat, binomial, and GKP codes.\cite{DuttonSEFI2026,DuttonElectronField2026} These frameworks provide strong protection against loss and dephasing but
do not explicitly incorporate the geometric constraints that govern photonic mode evolution. The SEFI field
architecture\cite{DuttonSEFI2026} and the Geometric Worldline Field Model (GWFM)\cite{DuttonElectronField2026} introduce stability surfaces and worldline
geometry as physically meaningful structures for describing photonic modes, errors, and correction. In this
work, Dynamic Entity Field Integration (DEFI) is introduced as the unifying field law that binds SEFI identity
and GWFM geometry into a single stability surface. This DEFI-centric framework provides a geometric
foundation for photonic QEC, enabling curvature-based stabilization, geometric filtering, and sovereignty
restoration across diverse photonic platforms. Figure~\ref{fig:stability_surface} illustrates the DEFI stability
surface that anchors this unified geometric interpretation.

Quantum error correction (QEC) in photonic systems
has traditionally been approached through algebraic or
phase-space constructions.\cite{DuttonSEFI2026} Cat codes stabilize coherentstate superpositions through engineered nonlinearities,
binomial codes use discrete photon-number symmetries,
and GKP codes embed logical states in a lattice of phasespace peaks. These methods have demonstrated strong
protection against loss and dephasing, yet they do not explicitly incorporate the geometric constraints that govern
photonic mode evolution. As photonic platforms continue to scale, the absence of a geometric foundation becomes increasingly limiting.

Photonic modes propagate along trajectories determined by boundary conditions, dispersion relations, and
field constraints.\cite{DuttonElectronField2026} These trajectories can be viewed as
worldlines in an abstract configuration space. In ideal
conditions, a mode follows a stable, physically permissible path. Perturbations such as scattering, absorption,
spectral drift, or mode mixing cause deviations from this
ideal trajectory. Traditional QEC frameworks describe
these deviations algebraically, but they do not provide
a geometric interpretation of stability or error displacement. Figure~\ref{fig:mode_evolution} provides an example
of how photonic worldlines evolve across a stability surface.

The SEFI field architecture introduces stability surfaces as geometric objects that encode the permissible
evolution of a field entity.\cite{DuttonSEFI2026} These surfaces arise from
sovereignty conditions that determine how a mode maintains its identity under ideal evolution. The Geometric
Worldline Field Model (GWFM)\cite{DuttonElectronField2026} extends this concept by
describing field entities as worldlines constrained to these
surfaces. Together, SEFI and GWFM provide a geometric language for understanding how photonic modes
evolve, how they deviate, and how they can be restored.
Figure~\ref{fig:sefi_invariants} shows the SEFI identity-layer
invariant surface that constrains permissible worldline identity.

In this work, Dynamic Entity Field Integration (DEFI)
is introduced as the unifying field law that binds SEFI
identity and GWFM geometry into a single stability surface.\cite{DuttonSEFI2026,DuttonElectronField2026} SEFI defines the identity-layer sovereignty conditions, GWFM defines the geometry-layer worldline structure, and DEFI provides the dynamical invariance that
fuses identity and geometry into unity. Stability surfaces
become DEFI stability manifolds on which sovereigntypreserving worldlines must lie, and deviations from these
manifolds acquire a precise geometric and dynamical interpretation. Figure~\ref{fig:error_displacement} illustrates
how photonic errors appear as geometric displacements above
the DEFI stability surface.

This geometric perspective is motivated by several considerations. First, photonic systems naturally exhibit
geometric structure through waveguide curvature, cavity boundary conditions, dispersion surfaces, and mode
profiles.\cite{DuttonElectronField2026} Second, geometric descriptions offer intuitive
interpretations of error processes: loss corresponds to inward displacement, dephasing corresponds to torsional
deviation, and mode mixing corresponds to curvatureinduced divergence. Third, geometric correction operations can be implemented through optical transformations that project displaced worldlines back onto their
DEFI stability surfaces, as shown in Figure~\ref{fig:correction_dynamics}.

By framing photonic QEC in terms of DEFI
stability surfaces and worldline geometry, the
SEFI/GWFM/DEFI approach provides a physically motivated alternative to algebraic bosonic codes.\cite{DuttonSEFI2026}
It complements existing methods while offering new
mechanisms for stabilizing photonic modes through curvature engineering, geometric filtering, and sovereignty
restoration. This section establishes the motivation
for a DEFI-centric geometric QEC framework and sets
the stage for the formal development presented in the
following sections.
\section{SEFI STABILITY GEOMETRY}

The SEFI field architecture provides the geometric
foundation for describing the stable evolution of photonic field entities.\cite{DuttonSEFI2026} In SEFI, a photonic mode is treated
as a single entity whose permissible evolution is constrained by sovereignty conditions—conditions that define how the entity maintains its identity under ideal,
lossless propagation. These sovereignty constraints give
rise to a stability surface, a geometric object on which
the mode’s worldline must lie. Figure~\ref{fig:sefi_invariants}
illustrates the SEFI identity-layer invariant surface.

In the DEFI-centric interpretation, this stability surface becomes a DEFI stability manifold.\cite{DuttonSEFI2026,DuttonElectronField2026}
SEFI supplies the identity-layer sovereignty conditions,
while DEFI supplies the dynamical invariance that ensures these conditions remain preserved along the worldline. The stability surface is therefore not merely a SEFI construct but
the identity-layer expression of DEFI’s unified field law.

Let $M$ denote the configuration space in which photonic worldlines evolve. A stability surface is a twodimensional manifold $\Sigma \subset M$ defined by a set of
sovereignty conditions,


\[
S(x) = 0,
\]


where $S$ encodes the geometric and physical constraints
characterizing ideal evolution.\cite{DuttonElectronField2026} These constraints may
arise from boundary conditions, dispersion relations, or
intrinsic field properties. The stability surface $\Sigma$ represents the set of points in $M$ where sovereignty conditions
are satisfied.

A worldline $\gamma(t)$ describing the evolution of a photonic
mode must satisfy


\[
\gamma(t) \in \Sigma
\]


under ideal conditions. The geometry of $\Sigma$ determines
the curvature and torsion of permissible trajectories. Let
$\kappa$ and $\tau$ denote the curvature and torsion of $\gamma(t)$, defined
through the Frenet–Serret frame.\cite{DuttonElectronField2026} Sovereignty requires
that these geometric quantities remain within bounds determined by the stability surface,


\[
\kappa = \kappa_{\Sigma}, \qquad \tau = \tau_{\Sigma},
\]


where $\kappa_{\Sigma}$ and $\tau_{\Sigma}$ are geometric invariants of the surface.

In the DEFI interpretation, these invariants become
identity-preserving dynamical constraints.\cite{DuttonSEFI2026}
DEFI requires that the worldline evolve in a manner that preserves SEFI identity invariants while remaining geometrically constrained by GWFM curvature and torsion
bounds. Thus, SEFI provides the identity manifold,
GWFM provides the geometry manifold, and DEFI binds
them into a single sovereignty-preserving stability structure.

The stability surface may be embedded in a higherdimensional configuration space, allowing photonic
modes to exhibit complex geometric behavior.\cite{DuttonElectronField2026}
For example, a saddle-shaped stability surface corresponds to a
mode whose permissible evolution exhibits both convergent and divergent curvature directions. Such surfaces
arise naturally in photonic systems with competing dispersion or boundary constraints. Figure~\ref{fig:stability_surface}
shows an example DEFI stability manifold.

The geometric structure of $\Sigma$ provides a mechanism for
identifying deviations from ideal evolution. If a worldline
$\gamma(t)$ leaves the stability surface, the sovereignty conditions are violated, and the deviation represents a photonic error. The magnitude and direction of this deviation can be quantified by the displacement vector


\[
d(t) = \gamma(t) - \Pi_{\Sigma}(\gamma(t)),
\]


where $\Pi_{\Sigma}$ denotes the geometric projection onto the stability surface.\cite{DuttonElectronField2026} The displacement vector encodes the geometric error and serves as the basis for correction operations. Figure~\ref{fig:error_displacement} depicts displacement above the stability surface.

In the DEFI-centric framework, this displacement vector represents a violation of sovereignty-preserving dynamical evolution.\cite{DuttonSEFI2026} Correction corresponds to restoring
the worldline to the DEFI stability manifold by removing
displacement and realigning curvature and torsion with
the geometric invariants of ideal evolution.

SEFI stability geometry therefore establishes the
identity-layer structure of photonic mode evolution.
When unified with GWFM geometry and DEFI dynamics, it forms the foundation of the geometric QEC framework developed in the following sections.
\section{GWFM WORLDLINE FORMALISM}

The Geometric Worldline Field Model (GWFM) describes field entities as worldlines evolving within a geometric configuration space.\cite{DuttonElectronField2026} In the context of photonic
quantum error correction, a photonic mode is represented
by a worldline $\gamma(t)$ whose evolution is constrained by
the stability surface $\Sigma$ introduced in Section 2. GWFM
provides the geometric machinery needed to characterize
ideal evolution, quantify deviations, and construct correction operations. Figure~\ref{fig:gwfm_geometry} illustrates
the curvature and torsion structure of worldline evolution.

In the DEFI-centric interpretation, GWFM becomes
the geometry-layer expression of the DEFI unified field
law.\cite{DuttonSEFI2026,DuttonElectronField2026} SEFI supplies the identity-layer sovereignty conditions, GWFM supplies the geometry-layer curvature
and torsion constraints, and DEFI binds these structures
into a single sovereignty-preserving dynamical evolution.
Thus, worldline geometry is not merely descriptive—it is
the geometric manifestation of DEFI’s stability requirements.

A worldline is defined as a smooth curve


\[
\gamma : \mathbb{R} \rightarrow M,
\]


where $M$ is the configuration space of photonic modes.
The geometric properties of $\gamma(t)$ are described using the
Frenet–Serret frame, consisting of the tangent, normal,
and binormal vectors,


\[
T(t),\; N(t),\; B(t),
\]


which satisfy the Frenet–Serret equations,\cite{DuttonElectronField2026}


\[
\frac{dT}{dt} = \kappa N,
\]




\[
\frac{dN}{dt} = -\kappa T + \tau B,
\]




\[
\frac{dB}{dt} = -\tau N.
\]


Here, $\kappa$ and $\tau$ denote the curvature and torsion of the
worldline. These quantities encode how the photonic
mode evolves geometrically and determine whether the
worldline remains on the stability surface.

Under ideal evolution, the worldline satisfies the
sovereignty conditions associated with the stability surface,


\[
\gamma(t) \in \Sigma,\qquad \kappa = \kappa_{\Sigma},\qquad \tau = \tau_{\Sigma}.
\]


The geometric invariants $\kappa_{\Sigma}$ and $\tau_{\Sigma}$ characterize the curvature and torsion of permissible trajectories on the stability surface. Any deviation from these invariants indicates a departure from ideal evolution.

In the DEFI interpretation, these geometric invariants
become dynamical invariants:\cite{DuttonSEFI2026} DEFI requires that worldline evolution preserve SEFI identity while remaining geometrically constrained by GWFM curvature and torsion
bounds. Thus, curvature and torsion are not merely geometric descriptors—they are DEFI stability conditions.

The GWFM formalism allows photonic modes to be
embedded in higher-dimensional configuration spaces,
enabling the description of complex geometric behavior
such as twisting, bending, or multi-mode coupling.\cite{DuttonElectronField2026}
For example, a photonic mode propagating through a curved
waveguide may exhibit nontrivial torsion, while a mode
experiencing dispersion-induced spreading may exhibit
curvature variations. These geometric features are naturally captured by the worldline formalism.

Errors arise when the worldline leaves the stability surface. Let $\Pi_{\Sigma}$ denote the geometric projection onto $\Sigma$.
The deviation of the worldline from the stability surface
is quantified by the displacement vector


\[
d(t) = \gamma(t) - \Pi_{\Sigma}(\gamma(t)).
\]


The magnitude of $d(t)$ represents the severity of the error, while its direction encodes the type of deviation. For
instance, displacement along the normal direction may
correspond to loss or amplitude damping, while displacement along the binormal direction may correspond to
dephasing or spectral drift.\cite{DuttonElectronField2026} Figure~\ref{fig:error_vector_field}
illustrates the error vector field associated with these deviations.

In the DEFI-centric framework, this displacement vector represents a violation of sovereignty-preserving dynamical evolution.\cite{DuttonSEFI2026} Correction corresponds to restoring
the worldline to the DEFI stability manifold by removing
displacement and realigning curvature and torsion with
the geometric invariants of ideal evolution.

The GWFM worldline formalism therefore provides
the geometry-layer structure of photonic mode evolution.
When unified with SEFI identity and DEFI dynamics, it
forms the foundation for the geometric QEC framework
developed in the following sections.
\section{ERROR DISPLACEMENT GEOMETRY}

Errors in photonic quantum systems manifest as deviations of the worldline $\gamma(t)$ from the stability surface $\Sigma$.
In the SEFI/GWFM framework, these deviations are described geometrically through displacement vectors that
quantify both the magnitude and direction of the error.\cite{DuttonElectronField2026}
In the DEFI-centric interpretation, error displacement
corresponds to a violation of sovereignty-preserving dynamical evolution: the worldline has left the DEFI stability manifold and must be restored. Figure~\ref{fig:error_displacement}
illustrates displacement above the DEFI stability surface.

Let $\Pi_{\Sigma}$ denote the geometric projection onto the stability surface. The displacement vector


\[
d(t) = \gamma(t) - \Pi_{\Sigma}(\gamma(t))
\]


captures the deviation of the worldline from ideal evolution.\cite{DuttonElectronField2026} The magnitude $\|d(t)\|$ represents the severity of
the error, while the direction of $d(t)$ encodes the type
of deviation. Because the stability surface is a geometric object embedded in the configuration space $M$, the
displacement vector naturally decomposes into components along the normal and binormal directions of the
Frenet–Serret frame.

Let $N(t)$ and $B(t)$ denote the normal and binormal
vectors of the worldline. The displacement vector can be
expressed as


\[
d(t) = d_{N}(t) N(t) + d_{B}(t) B(t),
\]


where $d_{N}(t)$ and $d_{B}(t)$ are scalar coefficients representing the normal and binormal components of the error.\cite{DuttonElectronField2026}
These components correspond to physically distinct error processes. Normal displacement typically represents
amplitude damping or photon loss, as it reflects inward
deviation toward regions of reduced field sovereignty. Binormal displacement corresponds to phase or spectral
drift, as it reflects torsional deviation from the stability
surface. Figure~\ref{fig:error_vector_field} shows the error vector field
associated with these deviations.

Curvature and torsion deviations provide additional insight into error formation. Deviations in curvature,


\[
\Delta\kappa(t) = \kappa(t) - \kappa_{\Sigma},
\]


indicate bending away from the stability surface, while
deviations in torsion,


\[
\Delta\tau(t) = \tau(t) - \tau_{\Sigma},
\]


indicate twisting or rotational drift.\cite{DuttonElectronField2026}
These geometric deviations complement the displacement vector and offer a
more complete description of photonic errors.

In the DEFI interpretation, these deviations correspond to violations of sovereignty-preserving dynamical
evolution.\cite{DuttonSEFI2026} DEFI requires that identity-layer invariants
(SEFI), geometry-layer invariants (GWFM), and dynamical invariants all remain synchronized. When curvature,
torsion, or displacement deviate from their DEFI stability values, the worldline has left the unified stability
manifold $S_{\mathrm{unified}}$.

The stability surface itself determines how errors accumulate. For example, a saddle-shaped stability surface
may exhibit regions where small deviations grow rapidly,
corresponding to photonic modes that are highly sensitive to perturbations. Conversely, surfaces with gentle
curvature may exhibit greater robustness.\cite{DuttonElectronField2026}
The geometric structure of $\Sigma$ thus plays a central role in determining
the error landscape of the photonic system.

Error displacement geometry provides a unified description of photonic errors that is independent of specific physical mechanisms.\cite{DuttonElectronField2026}
Whether the error arises
from scattering, absorption, dispersion, or mode mixing,
its geometric signature is captured by the displacement
vector and curvature/torsion deviations. In the DEFIcentric framework, these deviations represent violations
of sovereignty-preserving evolution and form the basis for
the correction operations developed in Section 5.
\section{CORRECTION MAPPING}

Correction operations in the SEFI/GWFM framework
restore sovereignty by mapping displaced worldlines back
onto the stability surface $\Sigma$.\cite{DuttonSEFI2026,DuttonElectronField2026}
Because errors are represented geometrically as displacement vectors, correction
becomes a projection process: the worldline is returned
to the nearest point on the stability surface, and its curvature and torsion are realigned with the geometric invariants of ideal evolution. Figure~\ref{fig:correction_dynamics}
illustrates sovereignty restoration through DEFI projection.

In the DEFI-centric interpretation, correction corresponds to restoring the worldline to the DEFI stability manifold.\cite{DuttonSEFI2026} SEFI provides the identity-layer invariants
that must be preserved, GWFM provides the geometrylayer curvature and torsion constraints, and DEFI provides the dynamical invariance that ensures sovereigntypreserving evolution. Thus, correction is not merely geometric—it is the restoration of DEFI unity.

Let $\gamma(t)$ denote the displaced worldline and let $\Pi_{\Sigma}$ denote the geometric projection onto the stability surface.
The corrected worldline is defined as


\[
\gamma_{\mathrm{corr}}(t) = \Pi_{\Sigma}(\gamma(t)).
\]


This projection removes the displacement vector,


\[
d(t) = \gamma(t) - \Pi_{\Sigma}(\gamma(t)),
\]


and restores the worldline to the stability surface.\cite{DuttonElectronField2026}
The projection operator $\Pi_{\Sigma}$ is determined by the geometry of
$\Sigma$ and may be implemented through optical transformations that selectively filter or redirect photonic modes.
Figure~\ref{fig:projection_mapping} shows the geometric projection map.

Correction requires not only returning the worldline
to the stability surface but also restoring its geometric
invariants. Let $\kappa(t)$ and $\tau(t)$ denote the curvature and
torsion of the displaced worldline. The corrected worldline must satisfy


\[
\kappa_{\mathrm{corr}}(t) = \kappa_{\Sigma}, \qquad \tau_{\mathrm{corr}}(t) = \tau_{\Sigma},
\]


where $\kappa_{\Sigma}$ and $\tau_{\Sigma}$ are the curvature and torsion of ideal
evolution.\cite{DuttonElectronField2026} These invariants may be restored through
optical phase shifts, dispersion compensation, or modeselective transformations that adjust the geometric properties of the photonic mode.

The correction mapping can be decomposed into normal and binormal components. Let $d_{N}(t)$ and $d_{B}(t)$ denote the normal and binormal components of the displacement vector. The correction operator acts as


\[
\gamma_{\mathrm{corr}}(t) = \gamma(t) - d_{N}(t) N(t) - d_{B}(t) B(t),
\]


removing both amplitude-related and phase-related deviations.\cite{DuttonElectronField2026} This geometric decomposition provides a unified
description of correction operations that is independent
of specific physical mechanisms.

In the DEFI interpretation, this mapping corresponds
to restoring sovereignty-preserving dynamical evolution.\cite{DuttonSEFI2026}
The worldline is returned to the DEFI stability manifold,
identity invariants are preserved, and geometric invariants are realigned. Correction therefore restores unity
across SEFI identity, GWFM geometry, and DEFI dynamics.

In photonic systems, correction may be implemented
through a variety of physical processes.\cite{DuttonElectronField2026}
Normal displacement may be corrected through amplitude stabilization or loss-compensating transformations, while binormal displacement may be corrected through phase stabilization or spectral filtering. The geometric nature of the
correction mapping ensures that these operations restore
sovereignty by returning the worldline to the stability
surface and realigning its geometric invariants.

The correction mapping developed in this section provides a physically motivated mechanism for restoring
photonic modes to ideal evolution. It forms the basis
for the photonic implementation strategies discussed in
Section 6.
\section{PHOTONIC IMPLEMENTATION}

The geometric framework developed in the preceding sections provides a physically motivated description
of photonic quantum error correction.\cite{DuttonSEFI2026,DuttonElectronField2026}
To apply this
framework in practical systems, the stability surface $\Sigma$,
worldline geometry, and correction mapping must be interpreted in terms of photonic modes, optical components, and experimentally accessible transformations. In
the DEFI-centric interpretation, photonic implementation corresponds to engineering systems that maintain
worldlines on the DEFI stability manifold through identity preservation (SEFI), geometric constraint (GWFM),
and sovereignty-preserving dynamics (DEFI). Figure~\ref{fig:qec_stack}
illustrates the SEFI\_QEC\_Stack architecture used to implement
these geometric operations.

Photonic modes propagate through physical structures
such as waveguides, fibers, cavities, and free-space optical
paths.\cite{DuttonElectronField2026} Each structure imposes geometric constraints on
the mode’s evolution, including boundary conditions, dispersion relations, and mode profiles. These constraints
define the stability surface $\Sigma$ for the photonic system. For
example, a waveguide with fixed transverse confinement
defines a stability surface corresponding to the permissible spatial mode profile, while a cavity with fixed resonance conditions defines a stability surface corresponding
to the permissible spectral structure.

In the DEFI interpretation, these stability surfaces become DEFI stability manifolds: geometric objects on
which sovereignty-preserving worldlines must lie.\cite{DuttonSEFI2026}
SEFI
provides the identity-layer invariants that define the permissible evolution, GWFM provides the geometry-layer
curvature and torsion constraints, and DEFI binds these
structures into a single unified stability requirement.

The worldline $\gamma(t)$ represents the evolution of the photonic mode within the configuration space $M$. In spatially structured systems, $\gamma(t)$ may correspond to the
trajectory of the mode’s centroid or phase front. In spectrally structured systems, $\gamma(t)$ may correspond to the
evolution of the mode’s frequency distribution. The curvature and torsion of the worldline encode how the mode
evolves under the influence of dispersion, nonlinearities,
or external perturbations.\cite{DuttonElectronField2026} Figure~\ref{fig:mode_evolution}
illustrates a photonic mode evolving across a stability surface.

Errors arise when the photonic mode deviates from its
stability surface.\cite{DuttonElectronField2026} Loss corresponds to inward displacement toward regions of reduced field sovereignty, while
dephasing corresponds to torsional deviation from the
stability surface. Mode mixing corresponds to curvatureinduced divergence, and spectral drift corresponds to binormal displacement. These geometric signatures provide a unified description of photonic errors that is independent of specific physical mechanisms. In the DEFIcentric interpretation, these errors correspond to violations of sovereignty-preserving dynamical evolution.

Correction operations may be implemented through
optical transformations that project displaced worldlines
back onto the stability surface.\cite{DuttonElectronField2026} Spatial corrections may
be achieved through mode-selective filtering, adaptive
optics, or waveguide engineering. Spectral corrections
may be achieved through dispersion compensation, phase
stabilization, or frequency-selective filtering. In each
case, the physical transformation acts as a geometric projection that removes the displacement vector and restores
the worldline to the stability surface.

In the DEFI interpretation, these correction operations
restore unity across SEFI identity, GWFM geometry, and
DEFI dynamics.\cite{DuttonSEFI2026} The worldline is returned to the DEFI
stability manifold, identity invariants are preserved, and
geometric invariants are realigned. Photonic implementation therefore becomes the engineering realization of
DEFI sovereignty restoration.

The SEFI/GWFM/DEFI framework is compatible
with a wide range of photonic platforms.\cite{DuttonElectronField2026}
In spontaneous
parametric down-conversion (SPDC) systems, stability
surfaces may correspond to the joint spectral amplitude
or spatial mode structure of the generated photons. In
integrated photonic circuits, stability surfaces may correspond to the guided mode profiles or resonance conditions of microresonators. In fiber-based systems, stability
surfaces may correspond to the permissible propagation
modes or dispersion surfaces. The geometric nature of
the framework ensures that it can be applied across diverse photonic architectures.

Photonic implementation of geometric correction operations provides a pathway for stabilizing quantum
information in continuous-variable systems.\cite{DuttonSEFI2026}
By interpreting photonic errors as geometric displacements
and correction operations as geometric projections, the
SEFI/GWFM/DEFI framework offers a physically motivated alternative to algebraic bosonic codes. It enables new mechanisms for stabilizing photonic modes
through curvature engineering, geometric filtering, and
sovereignty restoration.

The photonic implementation strategies developed in
this section form the basis for the numerical validation
presented in Section 7.
\section{VALIDATION}

To demonstrate the effectiveness of the geometric
correction framework developed in the preceding sections, we present a simple numerical example illustrating
stability deviation, error displacement, and sovereignty
restoration.\cite{DuttonElectronField2026,DuttonSEFI2026}
The purpose of this validation is not to
model a specific photonic device, but to show how geometric worldline analysis provides a clear and reproducible mechanism for identifying and correcting photonic errors. In the DEFI-centric interpretation, this example illustrates how worldlines depart from and are restored to the DEFI stability manifold. Figure~\ref{fig:validation_surface}
illustrates the saddle-shaped stability surface used in this example.

Consider a stability surface $\Sigma$ embedded in a threedimensional configuration space $M$. For validation, we
choose a saddle-shaped surface defined by


\[
\Sigma = \{(x, y, z) \in \mathbb{R}^{3} \mid z = \alpha x^{2} - \beta y^{2}\},
\]


where $\alpha$ and $\beta$ are positive constants.\cite{DuttonElectronField2026}
This surface exhibits both convergent and divergent curvature directions, providing a nontrivial geometric structure for testing stability and correction. The geometric invariants
of ideal evolution are determined by the curvature and
torsion of worldlines constrained to $\Sigma$.

An ideal worldline $\gamma_{\Sigma}(t)$ on the stability surface may
be parameterized as


\[
\gamma_{\Sigma}(t) = (t, \sin t, \alpha t^{2} - \beta \sin^{2} t),
\]


which satisfies the sovereignty conditions associated with
$\Sigma$.\cite{DuttonSEFI2026} The curvature and torsion of this worldline, denoted
$\kappa_{\Sigma}(t)$ and $\tau_{\Sigma}(t)$, serve as the geometric invariants of ideal
evolution. In the DEFI interpretation, these invariants
represent identity-preserving dynamical constraints.

To simulate an error, we introduce a displaced worldline


\[
\gamma(t) = \gamma_{\Sigma}(t) + \epsilon v(t),
\]


where $\epsilon$ is a small perturbation parameter and $v(t)$ is a
deviation vector chosen to represent a physically meaningful photonic error.\cite{DuttonElectronField2026}
For validation, we select


\[
v(t) = (0, 0, \cos t),
\]


corresponding to a binormal displacement that models
spectral or phase drift. Figure~\ref{fig:validation_displacement}
illustrates the displaced worldline oscillating above and below the stability surface.

The displaced worldline no longer
satisfies the sovereignty conditions, and its curvature and
torsion deviate from the geometric invariants of ideal evolution.

The displacement vector is computed as


\[
d(t) = \gamma(t) - \Pi_{\Sigma}(\gamma(t)),
\]


where $\Pi_{\Sigma}$ denotes the geometric projection onto the stability surface.\cite{DuttonElectronField2026}
For the saddle surface defined above, the
projection is given by


\[
\Pi_{\Sigma}(x, y, z) = (x, y, \alpha x^{2} - \beta y^{2}).
\]


The magnitude of the displacement vector,


\[
\|d(t)\| = |z - (\alpha x^{2} - \beta y^{2})|,
\]


quantifies the severity of the error.\cite{DuttonElectronField2026}
In this example, the
displacement magnitude oscillates with $\cos t$, reflecting
the periodic nature of the simulated spectral drift.

Correction is performed by applying the geometric projection,


\[
\gamma_{\mathrm{corr}}(t) = \Pi_{\Sigma}(\gamma(t)),
\]


which removes the displacement vector and restores the
worldline to the stability surface.\cite{DuttonSEFI2026}
After correction, the
curvature and torsion of the worldline satisfy


\[
\kappa_{\mathrm{corr}}(t) = \kappa_{\Sigma}(t), \qquad \tau_{\mathrm{corr}}(t) = \tau_{\Sigma}(t),
\]


demonstrating sovereignty restoration. Figure~\ref{fig:validation_correction}
shows the corrected worldline lying precisely on the stability surface.

In the DEFI-centric interpretation, this correction corresponds to restoring the worldline to the DEFI stability manifold.\cite{DuttonSEFI2026}
The identity-layer invariants (SEFI),
geometry-layer invariants (GWFM), and dynamical invariants (DEFI) are all realigned, demonstrating unity
restoration.

This validation example confirms that the
SEFI/GWFM/DEFI geometric framework provides
a clear and reproducible method for identifying and
correcting photonic errors.\cite{DuttonElectronField2026}
The geometric representation of stability deviation, error displacement, and
sovereignty restoration forms a robust foundation
for photonic quantum error correction across diverse
physical platforms.
\section{DISCUSSION}

The geometric framework developed in this work provides a physically motivated foundation for photonic
quantum error correction.\cite{DuttonSEFI2026,DuttonElectronField2026}
By representing photonic
modes as worldlines constrained to stability surfaces, the
SEFI/GWFM/DEFI approach offers a unified description of stability, error formation, and correction that is
independent of specific physical mechanisms. This geometric perspective complements existing bosonic codes
and suggests new pathways for stabilizing quantum information in photonic systems. Figure~\ref{fig:discussion_geometry}
illustrates the unified geometric structure underlying SEFI, GWFM, and DEFI.

One of the central advantages of the geometric approach is its ability to describe photonic errors in terms
of displacement vectors and curvature/torsion deviations.\cite{DuttonElectronField2026}
Traditional QEC frameworks typically classify errors algebraically, relying on operator-based descriptions
or phase-space symmetries. In contrast, the geometric
framework interprets errors as deviations from physically
meaningful stability surfaces. This interpretation provides intuitive insight into how errors arise and how they
can be corrected. For example, amplitude damping corresponds to inward normal displacement, while spectral
drift corresponds to binormal deviation. These geometric
signatures offer a natural classification of photonic errors
that is directly tied to the physical structure of the system.

The correction mapping developed in Section 5 demonstrates how geometric projection can restore sovereignty
by returning displaced worldlines to the stability surface.\cite{DuttonSEFI2026}
This projection-based correction mechanism is compatible with a wide range of photonic architectures, including waveguides, cavities, fibers, and integrated photonic
circuits. Because the correction mapping acts on geometric invariants rather than specific physical parameters, it
provides a flexible and platform-independent method for
stabilizing photonic modes. Figure~\ref{fig:projection_mapping}
illustrates this geometric correction mechanism.

The validation example presented in Section 7 illustrates the effectiveness of geometric projection in restoring sovereignty-preserving evolution.\cite{DuttonElectronField2026}
The saddle-shaped
stability surface used in the example demonstrates how
geometric features such as curvature and torsion influence error formation and correction. The displaced worldline oscillates above and below the stability surface due
to binormal perturbation, while the corrected worldline
returns precisely to the surface, demonstrating the effectiveness of geometric projection as a correction mechanism. Figure~\ref{fig:validation_correction} shows the corrected worldline
lying on the stability surface.

The geometric QEC framework developed here suggests several avenues for future research. First, stability
surfaces may be engineered to enhance robustness against
specific types of photonic errors. For example, surfaces
with gentle curvature may provide greater tolerance to
spectral drift, while surfaces with strong curvature may
enhance confinement against amplitude damping. Second, geometric filtering techniques may be developed to
remove displacement components selectively, enabling
targeted correction of specific error types. Third, the
SEFI/GWFM/DEFI framework may be extended to
multi-mode photonic systems, where stability surfaces
become higher-dimensional manifolds and worldlines
become multi-component geometric objects.

Finally, the geometric perspective may provide new insight into the relationship between photonic QEC and
other geometric frameworks in physics.\cite{DuttonElectronField2026}
For example, the
interpretation of photonic modes as worldlines on stability surfaces parallels geometric approaches in classical
mechanics, general relativity, and field theory. The unified SEFI/GWFM/DEFI framework may therefore serve
as a bridge between photonic QEC and broader geometric theories of physical systems.

In summary, the geometric QEC framework developed
in this work provides a physically motivated, platformindependent method for stabilizing photonic quantum
information. By interpreting photonic modes as worldlines constrained to stability surfaces and errors as geometric displacements, the SEFI/GWFM/DEFI approach
offers a unified description of stability, error formation,
and correction. This geometric perspective complements
existing bosonic codes and suggests new pathways for
stabilizing quantum information in photonic systems.
% ============================
% FIGURES
% ============================

\begin{figure}[p]
\centering
\includegraphics[width=0.48\textwidth]{fig9.jpg}
\caption{DEFI stability surface illustrating sovereignty-preserving evolution.}
\label{fig:stability_surface}
\end{figure}

\begin{figure}[p]
\centering
\includegraphics[width=0.48\textwidth]{fig8.jpg}
\caption{Photonic mode worldline evolving across a stability surface.}
\label{fig:mode_evolution}
\end{figure}

\begin{figure}[p]
\centering
\includegraphics[width=0.48\textwidth]{fig2.jpg}
\caption{SEFI identity-layer invariant surface defining permissible worldline identity.}
\label{fig:sefi_invariants}
\end{figure}

\begin{figure}[p]
\centering
\includegraphics[width=0.48\textwidth]{fig7.jpg}
\caption{Error displacement above the DEFI stability surface showing normal and binormal components.}
\label{fig:error_displacement}
\end{figure}

\begin{figure}[p]
\centering
\includegraphics[width=0.48\textwidth]{fig1.jpg}
\caption{Geometric correction dynamics restoring worldlines to the DEFI stability manifold.}
\label{fig:correction_dynamics}
\label{fig:projection_mapping}
\end{figure}

\begin{figure}[p]
\centering

\includegraphics[width=0.48\textwidth]{fig6.jpg}
\caption{GWFM curvature and torsion structure governing photonic worldline evolution.}
\label{fig:gwfm_geometry}
\end{figure}

\begin{figure}[p]
\centering
\includegraphics[width=0.48\textwidth]{fig5.jpg}
\caption{Normal and binormal error vector components associated with worldline displacement.}
\label{fig:error_vector_field}
\end{figure}

\begin{figure}[p]
\centering
\includegraphics[width=0.48\textwidth]{fig3.jpg}
\caption{SEFI\_QEC\_Stack architecture implementing geometric photonic QEC and the unified SEFI/GWFM/DEFI structure.}
\label{fig:qec_stack}
\label{fig:discussion_geometry}
\end{figure}

\begin{figure}[p]
\centering
\includegraphics[width=0.48\textwidth]{fig9.jpg}
\caption{Saddle-shaped stability surface used in the numerical validation example.}
\label{fig:validation_surface}
\end{figure}

\begin{figure}[p]
\centering
\includegraphics[width=0.48\textwidth]{fig10.jpg}
\caption{Displaced worldline oscillating above and below the stability surface and corrected worldline restored to $\Sigma$.}
\label{fig:validation_displacement}
\label{fig:validation_correction}
\end{figure}

% ============================
% BIBLIOGRAPHY
% ============================

\bibliographystyle{aipnum4-2}
\bibliography{references}

@book{NielsenChuang2010,
  author       = {M. A. Nielsen and I. L. Chuang},
  title        = {Quantum Computation and Quantum Information},
  publisher    = {Cambridge University Press},
  year         = {2010}
}

@article{ObrienPhotonics2009,
  author       = {J. L. O'Brien and A. Furusawa and J. Vu\v{c}kovi\'{c}},
  title        = {Photonic quantum technologies},
  journal      = {Nature Photonics},
  volume       = {3},
  pages        = {687--695},
  year         = {2009}
}

@misc{DuttonElectronField2026,
  author       = {J. Dutton},
  title        = {Geometric Worldline Foundations of the Electron Field},
  note         = {JMP manuscript JMP26-AR-02470 (2026)}
}

@misc{DuttonSEFI2026,
  author       = {J. Dutton},
  title        = {SEFI: The Single Entity Field Interpretation --- Formal Mathematical Construction},
  note         = {JMP manuscript JMP26-AR-02501 (2026)}
}

@misc{SEFIQECStack2026,
  author       = {J. Dutton},
  title        = {SEFI\_QEC\_Stack},
  howpublished = {\url{https://github.com/JasonSEFIDEFI/SEFI}},
  note         = {Accessed 2026}
}


\end{document}

